\newpage
\phantomsection
\label{prf:metric-diameter-set-monotone}
\LRAProofFor{thm:metric-diameter-set-monotone}

\begin{remark*}[Return]
\hyperref[thm:metric-diameter-set-monotone]{Return to Theorem}
\end{remark*}

\begin{theorem*}[Diameter Distance Sets are Monotone]
Let \((X,d)\) be a metric space and let \(A,B\subseteq X\).  If
\(A\subseteq B\), then \(D(A)\subseteq D(B)\).
\end{theorem*}

\begin{proof}[Professional Standard Proof]
\LRAProofBodyStart
Assume \(A\subseteq B\).  Let \(r\in D(A)\).  By definition of \(D(A)\), there
exist \(a,b\in A\) such that \(r=d(a,b)\).  Since \(A\subseteq B\), both
\(a\) and \(b\) also belong to \(B\).  Hence the same equality \(r=d(a,b)\)
witnesses that \(r\in D(B)\).  Therefore \(D(A)\subseteq D(B)\).
\end{proof}

\begin{proof}[Detailed Learning Proof]
\LRAProofBodyStart
To prove \(D(A)\subseteq D(B)\), take an arbitrary element of \(D(A)\).  Call
it \(r\).  Membership in \(D(A)\) means that \(r\) is not an abstract number
with no source: it is realized as a distance between two points of \(A\).
Thus there are points \(a,b\in A\) such that
\[
  r=d(a,b).
\]
The hypothesis \(A\subseteq B\) says that every point of \(A\) is also a point
of \(B\).  Therefore \(a,b\in B\).  The same two points now show that
\[
  r=d(a,b)
\]
is a distance between two points of \(B\).  Hence \(r\in D(B)\).  Since every
element of \(D(A)\) lies in \(D(B)\), we have \(D(A)\subseteq D(B)\).
\end{proof}

\begin{remark*}[Proof structure]
Unpack membership in the diameter distance set, transport the two witnesses
across the inclusion \(A\subseteq B\), and repack the same distance as an
element of \(D(B)\).
\end{remark*}

\begin{dependencies}
\begin{itemize}
  \item \hyperref[def:metric-diameter-distance-set]{Definition: Diameter Distance Set}
\end{itemize}
\end{dependencies}

\clearpage
